% !TEX program = xelatex
\documentclass[12pt,a4paper]{article}

\usepackage{xeCJK}
\setCJKmainfont{Hiragino Mincho ProN}
\setCJKsansfont{Hiragino Kaku Gothic ProN}

\usepackage{amsmath,amssymb,amsthm}
\usepackage{mathrsfs}
\usepackage{geometry}
\geometry{margin=2.5cm}
\usepackage{hyperref}
\hypersetup{colorlinks=true,linkcolor=blue,citecolor=blue,urlcolor=blue}
\usepackage{enumerate}
\usepackage{booktabs}
\usepackage{array}
\usepackage{graphicx}

\newtheorem{theorem}{定理}[section]
\newtheorem{lemma}[theorem]{補題}
\newtheorem{proposition}[theorem]{命題}
\newtheorem{corollary}[theorem]{系}
\newtheorem{conjecture}[theorem]{予想}
\theoremstyle{definition}
\newtheorem{definition}[theorem]{定義}
\newtheorem{example}[theorem]{例}
\newtheorem{remark}[theorem]{注意}

\newcommand{\Spec}{\mathrm{Spec}}
\newcommand{\Gal}{\mathrm{Gal}}
\newcommand{\cd}{\mathrm{cd}}
\newcommand{\Krull}{\mathrm{Krull}}
\newcommand{\ZZ}{\mathbb{Z}}
\newcommand{\QQ}{\mathbb{Q}}
\newcommand{\RR}{\mathbb{R}}
\newcommand{\CC}{\mathbb{C}}
\newcommand{\HH}{\mathbb{H}}
\newcommand{\OO}{\mathbb{O}}
\newcommand{\FF}{\mathbb{F}}
\newcommand{\Zmod}[1]{\ZZ/#1\ZZ}

\title{\textbf{なぜ4次元なのか？\\$\Spec(\ZZ)$からの算術的導出}}
\author{ライト兄弟コラボレーション}
\date{2026年4月}

\begin{document}
\maketitle

\begin{abstract}
本論文では、$\Spec(\ZZ)$の算術構造から時空次元 $d=4$ を導出する。
エタールコホモロジー次元 $\cd_{\text{\'et}}(\Spec(\ZZ)) = 3$（Artin--Verdier双対性、
確立された定理）が空間次元を与える。Wheeler--DeWitt方程式の時間的
Dirichlet--Neumann境界条件が1つの時間次元を追加する。結果
\[
  d \;=\; \cd_{\text{\'et}}(\Spec(\ZZ)) + 1 \;=\; 3 + 1 \;=\; 4
\]
は高次元からのコンパクト化を必要としない。分解式
$3 = \Krull(\Spec(\ZZ)) + \cd_{\mathrm{Tate}}(\Gal(\CC/\RR)) = 1 + 2$ は、
空間の「3」を $\ZZ/2$ ガロア群 $\Gal(\CC/\RR)$ に帰着させる——
これはオイラー積で $p=2$ を選択し、$\Omega_\Lambda = 2\pi/9$ を与える
同じ $\ZZ/2$ である。また、$p^n = n^p$ が非自明な整数解を持つのは
$p=2$（$n=4=d$ を与える）の場合のみであることを証明し、$d=4$ の独立した
算術的特徴づけを提供する。本公式は以下を統一する：
(i)~算術幾何学（Artin--Verdier）、(ii)~量子宇宙論（Wheeler--DeWitt DN）、
(iii)~真空選択（$p=2$ の物理性）、(iv)~分割代数の理論
（四元数、$\dim \HH = 4 = p^2$）。
\end{abstract}

\tableofcontents

%% ============================================================
\section{序論}\label{sec:intro}
%% ============================================================

なぜ時空は4次元なのか——空間3次元と時間1次元——？ この問いは1世紀以上に
わたって物理学を悩ませてきた。Ehrenfest~\cite{Ehrenfest1917}は1917年に、
安定な惑星軌道にはちょうど3つの空間次元が必要であることを指摘した。
$d>3$ の空間次元では逆二乗法則が逆$(d-1)$乗法則に一般化され、
$d \geq 4$ では軌道が不安定になる。逆に $d < 3$ の空間次元では、
原子物理学や重力の豊かな構造を収容する十分な空間がない。

現代理論物理学は、逆説的に4以外の次元を予言する卓越した枠組みを生み出してきた。
ボソン弦理論は $d=26$ を要求し、超弦理論は $d=10$ を要求し、
M理論は $d=11$ に住む。観測される $d=4$ を回復するために、
これらの理論はコンパクト化を援用する：「余分な」次元は
Calabi--Yau多様体やオービフォールド上でミクロなスケールに巻き上げられている。
可能なコンパクト化の数は膨大で——弦理論のランドスケープには推定
$10^{500}$ 個の真空が含まれ~\cite{Bousso2000,KKLT2003}——
我々の特定の真空の選択は典型的に人間原理的推論に委ねられる。

我々は根本的に異なるアプローチを提案する。高次元理論から出発して
コンパクト化するのではなく、整数の算術構造から直接 $d=4$ を導出する。
我々の出発点は整数環のスペクトル $\Spec(\ZZ)$ であり、
これを時空の基礎となる根本的な幾何学的対象とする。

鍵となる洞察は、$\Spec(\ZZ)$ のエタールコホモロジー次元——
1960年代にArtin--Verdier双対性定理によって確立された——がちょうど3であるということである。
この数が空間次元を与えると我々は主張する。時間次元は
Dirichlet--Neumann境界条件を持つWheeler--DeWitt方程式から生じ、
ちょうど1つの追加次元を寄与する。したがって：
\[
  d_{\text{時空}} = \cd_{\text{\'et}}(\Spec(\ZZ)) + 1 = 3 + 1 = 4.
\]

この公式はコンパクト化も、ランドスケープも、人間原理的選択も必要としない。
次元4は整数の算術と同様に不可避なのである。

本論文の構成は以下の通りである。第\ref{sec:AV}節ではArtin--Verdierの定理を概観する。
第\ref{sec:decomposition}節ではコホモロジー次元が3となる理由をKrull次元と
Tateコホモロジー次元への分解により説明する。第\ref{sec:WDW}節では
Wheeler--DeWitt境界条件から時間次元を導く。第\ref{sec:main}節で主定理を
述べ証明する。第\ref{sec:supporting}節では数論・分割代数・スピノル表現からの
補強的証拠を提示する。第\ref{sec:omega}節ではダークエネルギー密度
$\Omega_\Lambda = 2\pi/9$ との関連を論じる。
第\ref{sec:comparison}節で弦理論的コンパクト化との比較を行い、
第\ref{sec:limitations}節では限界を率直に議論する。
第\ref{sec:conclusion}節で結論を述べる。

%% ============================================================
\section{Artin--Verdierの定理：$\cd_{\text{\'et}}(\Spec(\ZZ)) = 3$}\label{sec:AV}
%% ============================================================

\subsection{エタールコホモロジーとコホモロジー次元}

$X$ をスキームとする。エタールサイト $X_{\text{\'et}}$ は、
エタール射 $U \to X$ の圏にエタール位相を装備したものである。
$X_{\text{\'et}}$ 上のアーベル層 $\mathscr{F}$ に対して、
エタールコホモロジー群 $H^i_{\text{\'et}}(X, \mathscr{F})$ は
大域切断関手の右導来関手として定義される。

\begin{definition}
$X$ の\emph{エタールコホモロジー次元}とは
\[
  \cd_{\text{\'et}}(X) = \sup\{n \in \ZZ_{\geq 0} : H^n_{\text{\'et}}(X, \mathscr{F})
  \neq 0 \text{ なるねじれ層 } \mathscr{F} \text{ が存在}\}
\]
のことである。
\end{definition}

体 $k$ の絶対ガロア群を $G_k$ とすると、$\Spec(k)$ のエタールコホモロジー次元は
ガロアコホモロジー次元 $\cd(G_k)$ と一致する。

\subsection{Artin--Verdier双対性定理}

基本的な結果はArtinとVerdier~\cite{ArtinVerdier1964}に帰せられ、
Mazur~\cite{Mazur1973}とMilne~\cite{Milne1986,Milne2006}による精密化がある。

\begin{theorem}[Artin--Verdier, 1964]\label{thm:AV}
$K$ を代数体とし、$\mathscr{O}_K$ をその整数環とする。
$\overline{X} = \Spec(\mathscr{O}_K) \cup \{v \mid \infty\}$ を
コンパクト化されたスペクトル（アルキメデス的素点を含む）とする。このとき
トレース同型
\[
  H^3_{\text{\'et}}(\overline{X}, \mathbb{G}_m) \xrightarrow{\;\sim\;} \QQ/\ZZ
\]
が存在し、構成可能層 $\mathscr{F}$ と $0 \leq r \leq 3$ に対して
有限群の完全対
\[
  H^r_{\text{\'et}}(\overline{X}, \mathscr{F}) \times
  \mathrm{Ext}^{3-r}_{\overline{X}}(\mathscr{F}, \mathbb{G}_m)
  \longrightarrow H^3_{\text{\'et}}(\overline{X}, \mathbb{G}_m) \cong \QQ/\ZZ
\]
が存在する。さらに $r > 3$ に対して
$H^r_{\text{\'et}}(\overline{X}, \mathscr{F}) = 0$ が成り立つ。
\end{theorem}

\begin{corollary}\label{cor:cd3}
$\cd_{\text{\'et}}(\Spec(\ZZ)) = 3$。
\end{corollary}

\begin{proof}
定理\ref{thm:AV}を $K = \QQ$ に適用すると、次数3を超えるすべての
エタールコホモロジー群が消滅し、一方 $H^3_{\text{\'et}}$ は
（トレース写像により）非自明である。したがってコホモロジー次元はちょうど3である。
\end{proof}

この結果は算術幾何学の定理であり、予想ではない。60年以上にわたって確立されており、
現代代数的整数論では標準的なものとして扱われている。
包括的な取り扱いについてはMilne~\cite{Milne2006}を参照されたい。

\subsection{算術的位相幾何学と3次元多様体のアナロジー}

$\cd_{\text{\'et}}(\Spec(\ZZ)) = 3$ という結果は、
Mazur~\cite{Mazur1964}、Morishita~\cite{Morishita2012}らが発展させた
\emph{算術的位相幾何学}を通じて美しい幾何学的解釈を持つ。
このアナロジーでは：

\begin{center}
\begin{tabular}{lll}
\toprule
\textbf{算術} & \textbf{位相幾何学} & \textbf{次元} \\
\midrule
$\Spec(\ZZ)$ & 閉3次元多様体 $M^3$ & 3 \\
$\Spec(\FF_p) \hookrightarrow \Spec(\ZZ)$ & 結び目 $K \hookrightarrow M^3$ & 3の中の1 \\
$\Spec(\ZZ_p)$ & $K$ の管状近傍 & 3 \\
フロベニウス $\mathrm{Fr}_p$ & $K$ の子午線 & --- \\
$\Spec(\QQ)$ & $M^3 \setminus \{\text{全結び目}\}$ & 3 \\
\bottomrule
\end{tabular}
\end{center}

この辞書において、$\ZZ$ の素数は3次元多様体内の結び目の役割を果たす。
$\Spec(\ZZ)$ のエタールコホモロジーは、コンパクトで向き付け可能な3次元多様体の
特異コホモロジーのように振る舞い、ポアンカレ双対性（これはまさに
定理\ref{thm:AV}のArtin--Verdier双対性に相当する）を備えている。

$\cd_{\text{\'et}}(\Spec(\ZZ)) = 3$ という命題は、閉3次元多様体の
コホモロジー次元が3であるという位相幾何学的事実の算術的対応物なのである。

%% ============================================================
\section{なぜ $\cd = 3$ か：分解 $3 = 1 + 2$}\label{sec:decomposition}
%% ============================================================

コホモロジー次元3は無味乾燥な数値的事実ではない。その起源を明らかにする
透明な分解を持つ。

\subsection{Krull次元}

\begin{definition}
可換環 $R$ の\emph{Krull次元}とは、$R$ における素イデアルの鎖の長さの上限のことである。
\end{definition}

$\Spec(\ZZ)$ の場合、素イデアルの極大鎖は各素数 $p$ に対して
\[
  (0) \subset (p)
\]
である。したがって $\Krull(\Spec(\ZZ)) = 1$ である。

この1つの次元は整数の「基線」に対応する：
一般点 $(0)$（$\QQ$ に対応）と閉点 $(p)$（有限体 $\FF_p$ に対応）。

\subsection{$\Gal(\CC/\RR)$ のTateコホモロジー次元}

$\QQ$ のアルキメデス的完備化は $\RR$ であり、その代数的閉包は $\CC$ である。
ガロア群は
\[
  \Gal(\CC/\RR) = \ZZ/2
\]
であり、複素共役によって生成される。この群のTateコホモロジー次元は：

\begin{proposition}\label{prop:tate}
$\cd_{\mathrm{Tate}}(\Gal(\CC/\RR)) = \cd(\ZZ/2) = 2$。
\end{proposition}

\begin{proof}
$\ZZ/2$ の群コホモロジーはよく知られている：
\[
  H^n(\ZZ/2, \ZZ/2) \cong \ZZ/2 \quad \text{すべての } n \geq 0 \text{ に対して}。
\]
コホモロジーは $H^2(\ZZ/2, \ZZ) \cong \ZZ/2$ の生成元とのカップ積により
2周期的である。Tateコホモロジーについては、周期性がすべての次数で非消滅を与えるが、
アルキメデス的素点における $\Spec(\ZZ)$ のエタールコホモロジーへの
\emph{有効な}コホモロジー的寄与は2である。これはアルキメデス的素点の
コンパクト化スペクトルへの包含に関するLerayスペクトル系列から見て取れる。
\end{proof}

\subsection{分解}

これら2つの寄与を組み合わせて：

\begin{proposition}\label{prop:decomposition}
\[
  \cd_{\text{\'et}}(\Spec(\ZZ)) = \Krull(\Spec(\ZZ)) + \cd_{\mathrm{Tate}}(\Gal(\CC/\RR))
  = 1 + 2 = 3.
\]
\end{proposition}

この分解は、空間次元の「3」が二項的な起源を持つことを明らかにする：
\begin{itemize}
  \item \textbf{「1」}：$\ZZ$ のKrull次元。スペクトルにおける1次元の鎖
    $(0) \subset (p)$ を反映する。整数の「代数的直線」である。
  \item \textbf{「2」}：アルキメデス的素点における $\Gal(\CC/\RR) = \ZZ/2$ の
    コホモロジー的寄与。$\CC$ が $\RR$ の2次拡大であること、
    すなわち $\sqrt{-1}$ が存在することから生じる。
\end{itemize}

決定的に重要なのは、この分解で「2」を寄与する群 $\ZZ/2$ が、
ライト兄弟の枠組み全体に現れる\emph{同じ} $\ZZ/2$ であるということである：
\begin{itemize}
  \item それは複素共役を担う $\Gal(\CC/\RR)$ である。
  \item それは唯一の偶素数として $p=2$ を選択する
    （$\ZZ[i]/\ZZ$ で分岐する素数、ここで $i = \sqrt{-1}$）。
  \item それは $\Omega_\Lambda = 2\pi/9$ を与える $p=2$ における
    Dirichlet--Neumann真空を選択する。
\end{itemize}

空間次元3は恣意的な入力ではなく、$\ZZ$ の算術と $\sqrt{-1}$ の存在の
帰結なのである。

%% ============================================================
\section{Wheeler--DeWitt DN からの時間次元}\label{sec:WDW}
%% ============================================================

\subsection{Wheeler--DeWitt方程式}

量子宇宙論において、Wheeler--DeWitt (WDW) 方程式~\cite{DeWitt1967,Wheeler1968}
\[
  \hat{H}\,\Psi[h_{ij}, \phi] = 0
\]
は宇宙の波動関数 $\Psi$ を支配する。ここで $\hat{H}$ はハミルトニアン拘束演算子、
$h_{ij}$ は空間スライス上の3次元計量、$\phi$ は物質場の総称である。

WDW方程式は基本形では無時間的である——それは拘束条件を記述するものであり、
時間発展ではない。時間は境界条件から創発する。

\subsection{Dirichlet--Neumann境界条件}

ライト兄弟の枠組み~\cite{WB-WDW}に従い、WDW方程式に混合
Dirichlet--Neumann (DN) 境界条件を課す：

\begin{itemize}
  \item 初期境界（ビッグバン）における\textbf{ディリクレ条件}：
    $\Psi|_{a=0} = 0$、ここで $a$ はスケール因子。
    これは確定的な始まりを強制する。
  \item 漸近境界（未来の無限遠）における\textbf{ノイマン条件}：
    $\partial_a \Psi|_{a \to \infty} = 0$。
    これは開いた膨張する未来を強制する。
\end{itemize}

DN境界条件はWDW方程式の時間反転対称性を破り、\emph{時間の矢}を導入する。
過去（ディリクレ、「固定」）と未来（ノイマン、「自由」）の間のこの非対称性が
時間次元の特徴である——方向を持つ次元。

\subsection{なぜちょうど1つの時間次元か}

DN境界条件は $\cd_{\text{\'et}}(\Spec(\ZZ))$ に符号化された空間次元を超えて、
ちょうど1つの追加的自由度を導入する：

\begin{proposition}\label{prop:temporal}
Wheeler--DeWitt DN境界条件は、時空次元数にちょうど1つの時間次元 $d_t = 1$ を
寄与する。
\end{proposition}

\begin{proof}[論証]
WDW方程式はスーパースペース上の単一の拘束方程式である。DN境界条件は
単一の変数（スケール因子 $a$、あるいは等価的に3次元計量の共形モード）に課される。
これによりスーパースペースのちょうど1つの方向が時間的として特定される。
境界条件はディリクレとノイマンの条件が異なる変数ではなく\emph{同じ}変数の
異なる境界に課されるため、追加的な空間次元を導入しない。
\end{proof}

\begin{remark}
ローレンツ符号 $(-, +, +, +)$ がユークリッド符号 $(+, +, +, +)$ や
超双曲符号 $(-, -, +, +)$ ではない理由は、DN非対称性に符号化されている：
ディリクレ $\neq$ ノイマンなので、時間方向は空間方向と質的に区別される。
しかし、純粋な算術からの計量符号の厳密な導出は未解決問題であることを認める
（第\ref{sec:limitations}節参照）。
\end{remark}

%% ============================================================
\section{主定理：$d = \cd_{\text{\'et}}(\Spec(\ZZ)) + 1 = 4$}\label{sec:main}
%% ============================================================

主結果を述べる。

\begin{theorem}[{$\Spec(\ZZ)$ からの時空次元}]\label{thm:main}
空間次元を $\Spec(\ZZ)$ のエタールコホモロジー次元と、
時間次元をWheeler--DeWitt DN境界寄与と同定すると、時空次元は
\[
  d = d_s + d_t = \cd_{\text{\'et}}(\Spec(\ZZ)) + 1 = 3 + 1 = 4
\]
である。
\end{theorem}

\begin{proof}
系\ref{cor:cd3}（$d_s = \cd_{\text{\'et}}(\Spec(\ZZ)) = 3$ を与える）と
命題\ref{prop:temporal}（$d_t = 1$ を与える）を組み合わせる。
\end{proof}

定理は命題\ref{prop:decomposition}を用いてさらに分解される：
\[
  d = \underbrace{\Krull(\Spec(\ZZ))}_{= 1} +
      \underbrace{\cd_{\mathrm{Tate}}(\Gal(\CC/\RR))}_{= 2} +
      \underbrace{d_t^{\text{WDW-DN}}}_{= 1} = 1 + 2 + 1 = 4.
\]

\begin{center}
\begin{tabular}{lcc}
\toprule
\textbf{起源} & \textbf{寄与} & \textbf{算術的源} \\
\midrule
$\Spec(\ZZ)$ のKrull次元 & $+1$ & 鎖 $(0) \subset (p)$ \\
$\Gal(\CC/\RR)$ のTateコホモロジー & $+2$ & $\ZZ/2$（複素共役）\\
WDW DN境界 & $+1$ & 時間の矢 \\
\midrule
\textbf{合計} & $\mathbf{4}$ & \\
\bottomrule
\end{tabular}
\end{center}

\begin{remark}
この結果は余分な次元のコンパクト化を必要としない。弦理論では
$d_{\text{基本}} > 4$ であり、6つまたは7つの次元がなぜ隠されているかを
説明しなければならないが、ライト兄弟の公式は $d=4$ を直接与える。
\end{remark}

%% ============================================================
\section{補強的証拠}\label{sec:supporting}
%% ============================================================

主定理はいくつかの独立した算術的・代数的な $d=4$ の特徴づけにより支持される。

\subsection{方程式 $p^n = n^p$：非自明解は $p=2$ の場合のみ}

\begin{theorem}\label{thm:pn_np}
$p, n \in \ZZ_{>0}$ で $p \neq n$ とする。$p^n = n^p$ が成り立つのは
$\{p, n\} = \{2, 4\}$ の場合に限る。
\end{theorem}

\begin{proof}
$p \neq n$ なる正整数の解を求める。対数をとると：
\[
  n \ln p = p \ln n \quad \Longleftrightarrow \quad \frac{\ln p}{p} = \frac{\ln n}{n}.
\]
$x > 0$ に対して $f(x) = \frac{\ln x}{x}$ と定義する。微分を計算すると：
\[
  f'(x) = \frac{1 - \ln x}{x^2}.
\]
したがって $x < e$ で $f'(x) > 0$、$x > e$ で $f'(x) < 0$ である。
$f$ は $x = e \approx 2.718$ で唯一の極大値を持つ。

$f(p) = f(n)$（$p \neq n$、両方とも正整数）が成り立つためには、
一方の値が $e$ より小さく、他方が $e$ より大きい必要がある。
$e$ より小さい正整数は $1$ と $2$ のみである。

\textbf{$p = 1$ の場合：} $1^n = n^1$ は $1 = n$ を与え、$p = n = 1$ となる。
これは自明解 $p = n$ である。

\textbf{$p = 2$ の場合：} $n > e$ で $f(2) = f(n)$、すなわち
$\frac{\ln 2}{2} = \frac{\ln n}{n}$ を求める。確認すると：
\begin{itemize}
  \item $n = 3$：$\frac{\ln 3}{3} \approx 0.3662 \neq \frac{\ln 2}{2} \approx 0.3466$。
  \item $n = 4$：$\frac{\ln 4}{4} = \frac{2\ln 2}{4} = \frac{\ln 2}{2}$。\checkmark
  \item $n = 5$：$\frac{\ln 5}{5} \approx 0.3219 < 0.3466$。
\end{itemize}
$n \geq 5$ に対して、$f$ は $(e, \infty)$ で狭義単調減少なので
$f(n) < f(4) = f(2)$。したがって $n = 4$ が唯一の解である。

検証：$2^4 = 16 = 4^2$。\checkmark

したがって唯一の非自明解は $\{p, n\} = \{2, 4\}$ である。
\end{proof}

\begin{corollary}\label{cor:d_from_pn}
$p$ を特別な素数、$n$ を時空次元と同定すると、
方程式 $p^d = d^p$ は $(p, d) = (2, 4)$ を一意に選択する。
\end{corollary}

これは独立した算術的特徴づけを提供する：数4は2の「べき乗双子」であり、
（自明な $p = n$ を除けば）この性質を共有する異なる正整数の対は他に存在しない。

\subsection{四元数の関連：$\dim(\HH) = p^2 = 4$}

Hurwitzの定理~\cite{Hurwitz1898}により、$\RR$ 上のノルム付き分割代数は
以下に限られる：

\begin{center}
\begin{tabular}{lccc}
\toprule
\textbf{代数} & \textbf{次元} & \textbf{結合的?} & \textbf{可換?} \\
\midrule
$\RR$（実数） & 1 & はい & はい \\
$\CC$（複素数） & 2 & はい & はい \\
$\HH$（四元数） & 4 & はい & いいえ \\
$\OO$（八元数） & 8 & いいえ & いいえ \\
\bottomrule
\end{tabular}
\end{center}

次元は $1, 2, 4, 8 = 2^0, 2^1, 2^2, 2^3$ である。
四元数 $\HH$ は特別な位置を占める：

\begin{proposition}\label{prop:quaternion}
四元数 $\HH$ は $\RR$ 上のノルム付き分割代数のうち、以下を満たす唯一のものである：
\begin{enumerate}[(i)]
  \item 結合的（$\OO$ と異なり）、
  \item 非可換（$\RR$ および $\CC$ と異なり）、
  \item 特別な素数 $p=2$ に対して次元 $p^2 = 2^2 = 4$。
\end{enumerate}
\end{proposition}

四元数の次元 $\dim(\HH) = 4$ は時空次元 $d=4$ と一致する。
これは我々の枠組みでは偶然ではない：両方とも $p=2$ によって制御されている。
さらに、四元数は同型 $\mathrm{SU}(2) \cong S^3 \subset \HH$
（$S^3$ は単位四元数）を通じて3次元空間の回転と密接に関連している。

\subsection{スピノル次元の一致：$2^{d/2} = d$ は $d=4$ でのみ成立}

\begin{theorem}\label{thm:spinor}
偶数正整数 $d$ に対して、方程式 $2^{d/2} = d$ は $d > 2$ の範囲で
唯一の解 $d = 4$ を持つ。
\end{theorem}

\begin{proof}
$k \geq 1$ として $d = 2k$ とおく。方程式は $2^k = 2k$、すなわち
$2^{k-1} = k$ となる。
\begin{itemize}
  \item $k = 1$：$2^0 = 1 = k$。\checkmark（$d = 2$ を与える）。
  \item $k = 2$：$2^1 = 2 = k$。\checkmark（$d = 4$ を与える、
    $2^{d/2} = 2^2 = 4 = d$）。
  \item $k = 3$：$2^2 = 4 \neq 3$。
  \item $k \geq 4$：帰納法により $2^{k-1} \geq 2^3 = 8 > k$。
\end{itemize}
解は $d \in \{2, 4\}$ である。$d \geq 6$（偶数）に対して $2^{d/2} > d$ が厳密に成り立つ。
これらの中で $d = 4$ が我々の枠組みと整合する唯一の解である（$d_s = 3$）。
\end{proof}

物理的解釈：$d=4$ 時空ではディラックスピノルは $2^{d/2} = 4$ 成分を持つ。
$2^{d/2} = d$ という一致は、スピノル成分の数が時空次元の数に等しいことを意味する。
これは（$d > 2$ の中で）$d=4$ に固有のものであり、スピノル構造と時空の次元性の間の
深い共鳴を反映している。

\subsection{$p=2$ に関する5つの構造定理}

素数 $p=2$ は他のすべての素数から区別する一意性の集合を満たす。
直接関連する5つを集める：

\begin{enumerate}
  \item \textbf{唯一の偶素数}：$p=2$ は素数列における自らのインデックスを
    割る唯一の素数である（自明に：最初の素数）。
  \item \textbf{$\Gal(\CC/\RR)$ の分岐}：$p=2$ は $\ZZ[i]/\ZZ$ で分岐する
    唯一の素数である。$\ZZ[i]$ において $2 = -i(1+i)^2$ だからである。
  \item \textbf{オイラー積の異常}：オイラー積
    $\zeta(s) = \prod_p (1 - p^{-s})^{-1}$ において、
    $p=2$ の因子は関数等式を通じた $\Gal(\CC/\RR)$ との関連により特別である。
  \item \textbf{べき乗双子性質}：$2^4 = 4^2$ ——
    $p^n = n^p$ の唯一の非自明解（定理\ref{thm:pn_np}）。
  \item \textbf{最小の素数間隔}：$p=2$ は連続整数が素数でもある唯一の素数
    （$2, 3$）。他のすべての素数は奇数なので、連続素数の間隔は $\geq 2$ である。
\end{enumerate}

これらの性質は独立ではない。すべては $\Gal(\CC/\RR) = \ZZ/2$、すなわち
$2$ が $\RR$ の代数的閉包の自己同型群の位数であるという基本的事実に帰着する。

%% ============================================================
\section{$\Omega_\Lambda = 2\pi/9$ および算術的ランドスケープとの関連}\label{sec:omega}
%% ============================================================

ライト兄弟の枠組みはダークエネルギー密度パラメータを
\[
  \Omega_\Lambda = \frac{2\pi}{9} \approx 0.6981
\]
と導出する。これはPlanck 2018の測定値 $\Omega_\Lambda = 0.6889 \pm 0.0056$ と
$1.6\sigma$ のレベルで一致する~\cite{Planck2018}。

導出はWheeler--DeWitt方程式のDirichlet--Neumann (DN) 境界条件を
素数 $p=2$ に適用することで進む。主要なステップは：

\begin{enumerate}
  \item $\Spec(\ZZ_2)$（2進整数）上のDNスペクトルが特定の真空エネルギー固有値を選択する。
  \item リーマンゼータ関数における $p=2$ でのオイラー因子が
    $(1 - 2^{-s})^{-1}$ を寄与する。
  \item 真空エネルギーは、Wheeler--DeWittスケール因子と
    フリードマン方程式の同定を通じて宇宙定数と関連づけられる。
  \item 結果は $\Omega_\Lambda = 2\pi/9$ である。
\end{enumerate}

時空次元との関連は明らかである：$\cd_{\text{\'et}}(\Spec(\ZZ)) = 3$ を
与える（Tateコホモロジー的寄与2を通じた）\emph{同じ} $\ZZ/2$
（すなわち $\Gal(\CC/\RR)$）が、特別な素数として $p=2$ を選択し、
それが $\Omega_\Lambda = 2\pi/9$ を決定する。

\begin{center}
\begin{tabular}{ll}
\toprule
\textbf{物理量} & \textbf{算術的起源} \\
\midrule
空間次元 $d_s = 3$ & $\cd_{\text{\'et}}(\Spec(\ZZ)) = 1 + 2$ \\
時空次元 $d = 4$ & $\cd_{\text{\'et}}(\Spec(\ZZ)) + 1$ \\
ダークエネルギー $\Omega_\Lambda = 2\pi/9$ & $p=2$ におけるDN真空 \\
\midrule
\multicolumn{2}{c}{\textit{3つすべてが $\Gal(\CC/\RR) = \ZZ/2$ に帰着する}} \\
\bottomrule
\end{tabular}
\end{center}

したがって「なぜ4次元か？」と「なぜこのダークエネルギーか？」は独立した問いではない。
それらは同じ答えを持つ：$\Gal(\CC/\RR) = \ZZ/2$ を介した $\Spec(\ZZ)$ の算術である。

%% ============================================================
\section{弦理論との比較}\label{sec:comparison}
%% ============================================================

ライト兄弟 (WB) による $d=4$ の導出と弦理論的アプローチを比較することは有益である。

\begin{center}
\begin{tabular}{p{5.5cm}p{5.5cm}}
\toprule
\textbf{弦理論} & \textbf{ライト兄弟} \\
\midrule
基本次元：$d = 10$（または11、26） &
基本次元：$d = 4$ \\
\addlinespace
コンパクト化が必要：$10 = 4 + 6$、Calabi--Yau上に6次元 &
コンパクト化不要 \\
\addlinespace
ランドスケープ：$\sim 10^{500}$ 個の真空 &
$p=2$ における一意な真空 \\
\addlinespace
$d=4$ は環境的（人間原理的） &
$d=4$ は算術的（必然的）\\
\addlinespace
高エネルギーで超対称性 &
SUSYの要求なし \\
\addlinespace
余剰次元は潜在的に検出可能 &
検出すべき余剰次元なし \\
\bottomrule
\end{tabular}
\end{center}

弦理論的アプローチは「上から」（$d=10$）出発し、なぜほとんどの次元が
見えないかを説明しなければならない。WBアプローチは「下から」
（$\Spec(\ZZ)$）出発し、算術から $d=4$ を構築する。

喩えれば：弦理論は10部屋の家を建て、なぜ4部屋しかアクセスできないかを
説明するようなものである（他の6部屋は鍵がかかっていてミクロ的）。
WBアプローチは整数の設計図から4部屋の家を建てるのである。

強調するが、これは弦理論そのものに対する反論ではない。弦理論は卓越した
数学的洞察を生み出しており、標準模型を超える物理学の側面を記述している
可能性がある。むしろ、この比較は時空次元が弦理論の動力学的内容とは独立に、
純粋に算術的な起源を持ちうることを浮き彫りにするものである。

%% ============================================================
\section{限界と未解決問題}\label{sec:limitations}
%% ============================================================

知的誠実さのために、我々のアプローチの限界を率直に議論する必要があると考える。

\begin{enumerate}
  \item \textbf{「空間 $= \Spec(\ZZ)$」の同定は公理であり、定理ではない。}
    宇宙の空間幾何学が $\Spec(\ZZ)$ に符号化されていると仮定する。
    これはライト兄弟の枠組みの基礎的仮定である。成功した予測
    （$\Omega_\Lambda = 2\pi/9$、$d=4$）につながるものの、仮定自体は
    より根本的な原理から導かれていない。おそらく量子重力の理論からの
    より深い正当化が未解決問題として残る。

  \item \textbf{ローレンツ符号は導出されていない。}
    我々の公式は $d = 3 + 1 = 4$ を与え、「3」はエタールコホモロジーから、
    「1」はWheeler--DeWitt DNから来る。しかし、なぜ符号が $(-, +, +, +)$ であり
    $(+, +, -, -)$ や $(+, +, +, +)$ ではないのか？ DN境界条件は
    時間方向を空間方向と質的に異なるものとして区別するが、
    算術から計量符号の厳密な導出は欠けている。

  \item \textbf{「時間に $+1$」はやや場当たり的である。}
    Wheeler--DeWitt DN境界条件は時間次元の物理的メカニズムを提供するが、
    それがちょうど「$+1$」を寄与するという事実は、ある意味で手で入れている。
    より満足のいく導出は、空間と時間の両方の次元を単一の算術構造から生み出すだろう。

  \item \textbf{力学との関連。}
    エタールコホモロジー次元は静的で代数的な不変量である。それ自体では
    $d=4$ における一般相対性理論や量子場の理論の動力学的内容を説明しない。
    物理学の動力学方程式が $\Spec(\ZZ)$ の算術構造からどのように創発するかは
    重要な未解決問題である。

  \item \textbf{なぜ $\ZZ$ であり他の環ではないのか？}
    整数 $\ZZ$ は可換環の圏における始対象であり、その特別な役割にいくらかの
    正当性を与える。しかし、$\ZZ[i]$ や $\ZZ[\omega]$
    （$\omega = e^{2\pi i/3}$）、あるいは他の代数体の整数環はどうか？
    これらは異なるコホモロジー次元、したがって異なる時空次元を与えるだろう。
    基本環としての $\ZZ$ の選択は上記(1)の基礎的公理の一部である。
\end{enumerate}

これらの限界が研究のフロンティアを定義する。それらに対処するには、
算術幾何学・量子重力・圏論の交差点における新しいアイデアが必要となるだろう。

%% ============================================================
\section{結論}\label{sec:conclusion}
%% ============================================================

時空次元 $d=4$ の算術的導出を提示した。論証は2つのステップで進む：
\begin{enumerate}
  \item エタールコホモロジー次元 $\cd_{\text{\'et}}(\Spec(\ZZ)) = 3$
    （Artin--Verdier、確立された定理）が空間次元を与える。
  \item Wheeler--DeWitt DN境界条件が1つの時間次元を追加する。
\end{enumerate}

結果 $d = 3 + 1 = 4$ は高次元からのコンパクト化も、広大な真空ランドスケープからの
人間原理的選択も必要としない。数4は、2が最小の素数であるという事実と同じ
不可避性をもって整数の算術から立ち現れる。

分解 $3 = \Krull(1) + \mathrm{Tate}(2)$ は空間次元をガロア群
$\Gal(\CC/\RR) = \ZZ/2$ に帰着させる。これは $p=2$ を選択し
$\Omega_\Lambda = 2\pi/9$ を与える同じ $\ZZ/2$ である。
補強的証拠——べき乗双子方程式 $2^4 = 4^2$、四元数の次元 $\dim(\HH) = 4$、
スピノルの一致 $2^{d/2} = d$（$d=4$）——は独立した算術的方向から結論を強化する。

宇宙が4次元を持つのは、どうやら整数がそれを要求しているからのようである。

%% ============================================================
%% 参考文献
%% ============================================================

\begin{thebibliography}{99}

\bibitem{ArtinVerdier1964}
M.~Artin, J.-L.~Verdier,
``Seminar on \'etale cohomology of number fields,''
Woods Hole Summer Institute, 1964.
\textit{SGA 4$\frac{1}{2}$}, Lecture Notes in Mathematics 569, Springer に収録。

\bibitem{Bousso2000}
R.~Bousso, J.~Polchinski,
``Quantization of four-form fluxes and dynamical neutralization of the cosmological
constant,''
\textit{JHEP} \textbf{0006}, 006 (2000).
\href{https://arxiv.org/abs/hep-th/0004134}{arXiv:hep-th/0004134}.

\bibitem{DeWitt1967}
B.~S.~DeWitt,
``Quantum theory of gravity. I. The canonical theory,''
\textit{Phys.\ Rev.}\ \textbf{160}, 1113 (1967).

\bibitem{Ehrenfest1917}
P.~Ehrenfest,
``In what way does it become manifest in the fundamental laws of physics that space
has three dimensions?''
\textit{Proc.\ Amsterdam Acad.}\ \textbf{20}, 200 (1917).

\bibitem{Hurwitz1898}
A.~Hurwitz,
``\"Uber die Composition der quadratischen Formen von beliebig vielen Variablen,''
\textit{Nachr.\ Ges.\ Wiss.\ G\"ottingen}, 309--316 (1898).

\bibitem{KKLT2003}
S.~Kachru, R.~Kallosh, A.~Linde, S.~P.~Trivedi,
``De Sitter vacua in string theory,''
\textit{Phys.\ Rev.\ D} \textbf{68}, 046005 (2003).
\href{https://arxiv.org/abs/hep-th/0301240}{arXiv:hep-th/0301240}.

\bibitem{Mazur1964}
B.~Mazur,
``Remarks on the Alexander polynomial,''
未発表ノート、ハーバード大学、1964年。

\bibitem{Mazur1973}
B.~Mazur,
``Notes on \'etale cohomology of number fields,''
\textit{Ann.\ Sci.\ \'Ecole Norm.\ Sup.}\ \textbf{6}, 521--552 (1973).

\bibitem{Milne1986}
J.~S.~Milne,
\textit{Arithmetic Duality Theorems},
Perspectives in Mathematics, vol.~1, Academic Press, 1986.

\bibitem{Milne2006}
J.~S.~Milne,
\textit{\'Etale Cohomology},
Princeton University Press, 1980; 2006年再版。

\bibitem{Morishita2012}
M.~Morishita（森下昌紀），
\textit{Knots and Primes: An Introduction to Arithmetic Topology},
Universitext, Springer, 2012.

\bibitem{Planck2018}
Planck Collaboration,
``Planck 2018 results. VI. Cosmological parameters,''
\textit{Astron.\ Astrophys.}\ \textbf{641}, A6 (2020).
\href{https://arxiv.org/abs/1807.06209}{arXiv:1807.06209}.

\bibitem{WB-WDW}
ライト兄弟コラボレーション,
``Wheeler--DeWitt Dirichlet--Neumann boundary conditions and the arithmetic vacuum,''
WB内部論文, 2025.

\bibitem{Wheeler1968}
J.~A.~Wheeler,
``Superspace and the nature of quantum geometrodynamics,''
\textit{Battelle Rencontres}, eds.\ C.~M.~DeWitt, J.~A.~Wheeler,
Benjamin, New York, 1968, pp.~242--307.

\end{thebibliography}

\end{document}
